\documentclass[12pt, letterpaper]{article}

%-------Packages---------
\usepackage{amssymb,amsfonts}
\usepackage{mathptmx}
\usepackage{amsthm}
\usepackage{graphicx}
\usepackage[all,arc]{xy}
\usepackage{enumerate}
\usepackage{bbm}
\usepackage{dsfont}
\usepackage{mathrsfs}
\usepackage{tikz-cd}
\usepackage{setspace}
\usepackage{import}
\usepackage[utf8]{inputenc}
\usepackage{hyperref}
\usepackage{tikz}
\usepackage{float}
\usepackage[dvipsnames]{xcolor}
\usepackage{fancyhdr}
\usepackage[nointegrals]{wasysym}

\usepackage[left=1in,top=1in,right=1in,bottom=1in]{geometry}

\usepackage{mathtools}
\usepackage{setspace}
\usepackage{cleveref}
\usepackage{multicol}
\usepackage{mathpazo}
\usepackage{tcolorbox}
\tcbuselibrary{theorems,skins,breakable}
\usepackage{xparse}

\usepackage{xcolor, ifthen, tcolorbox}
\tcbuselibrary{theorems,skins,breakable}
% Theorem System
% The following environments are provided:
%   New Problem:        \begin{problem} ...
%   New Question Header:   \qh (then followed by subproblems)
%   Subproblem:     \begin{subproblem} ...
%   Problem Proof:  \begin{problemproof} ...
%   Proof:          \\begin{pf}{color} ...
%   Remark:         \rmk (sentence), \rmkb (block)

% Define custom colors
\definecolor{thmscol}{HTML}{F4565B} %For Problems
\definecolor{lemscol}{HTML}{FE844C} %For Lemmes
\definecolor{prpscol}{HTML}{00A7EF} %For proposition
\definecolor{corscol}{HTML}{23DAFF} %For corrolaries
\definecolor{rmkscol}{HTML}{6DEEC5} %For remarks
\definecolor{defscol}{HTML}{18C991} %For definitions
\definecolor{exmscol}{HTML}{7DB800} %For examples
\definecolor{clmscol}{HTML}{165c58} %For claims

%Change between dark(black)/light(white) mode
\newcommand{\colormode}{white}
\ifthenelse{\equal{\colormode}{white}}
      {\newcommand{\TextColor}{black}}
      {\newcommand{\TextColor}{white}}
\newcommand{\pbcolormain}{thmscol!80!\colormode}
\newcommand{\pbcolorsecond}{thmscol!38!\colormode}


% ============================
% Problem Counter
% ============================
\newcounter{subproblemcounter}
\newcounter{problemcounter}

\newcommand{\q}{
    \setcounter{subproblemcounter}{0}%
    \stepcounter{problemcounter} % Increment problem counter
    \color{\TextColor}Problem \arabic{problemcounter}: % Display the problem counter
}

\newcommand{\stepsub}{
    \stepcounter{subproblemcounter}%
    \color{\TextColor}(\alph{subproblemcounter})
}
% ============================
% Problem
% ============================
\newtcolorbox{pb}[1]{
  enhanced,
  frame hidden,
  titlerule=0mm,
  toptitle=1mm,
  bottomtitle=1mm,
  fonttitle=\bfseries\large,
  coltitle=black,
  colbacktitle=\pbcolormain,
  colback=\pbcolorsecond,
  title={\q #1},
  coltext=\TextColor
}

\newtcolorbox{spb}[1]{
  enhanced,
  frame hidden,
  titlerule=0mm,
  toptitle=1mm,
  bottomtitle=1mm,
  fonttitle=\bfseries\large,
  coltitle=black,
  colbacktitle=\pbcolormain,
  colback=\pbcolorsecond,
  title={\stepsub #1},
  coltext=\TextColor,
  attach boxed title to top left={xshift=3mm,yshift=-2mm},
  boxed title style={
    colback=\pbcolormain,
    colframe=\pbcolormain,
  }
}

\newtcolorbox{pbh}[1]{
    enhanced,
    frame hidden,
    titlerule=0mm,
    toptitle=1mm,
    bottomtitle=1mm,
    fonttitle=\bfseries\large,
    coltitle=black,
    colbacktitle=\pbcolormain,
    colback=\pbcolorsecond,
    title={\q #1},
    boxrule=0pt,
    interior hidden,
    attach boxed title to top left={xshift=3mm,yshift=-2mm},
    boxed title style={
        colback=\pbcolormain,
        colframe=\pbcolormain,
    }
}

\newenvironment{problem}[1]{%
  \newpage
  \begin{pb}{#1}
    }{
  \end{pb}
}

\newenvironment{subproblem}[1]{%
  \begin{spb}{#1}
    }{
  \end{spb}
}

\newenvironment{problemproof}{%
    \begin{pf}{\pbcolormain}
        }{
    \end{pf}
}

\newcommand{\qh}[1][]{
    \newpage
    \begin{pbh}{#1}\end{pbh} % Display the problem counter
}

% ============================
% Claim
% ============================

\newtcbtheorem[number within=problemcounter]{claim}{Claim}
{
    enhanced,
    frame hidden,
    titlerule=0mm,
    toptitle=1mm,
    bottomtitle=1mm,
    fonttitle=\bfseries\large,
    coltitle=black,
    colbacktitle=clmscol!50!\colormode,
    colback=clmscol!20!\colormode,
}{clm}

\NewDocumentCommand{\clm}{m+m}{
    \begin{claim*}{#1}{}
        #2
    \end{claim*}
}

\newenvironment{clmpf}{
	{\noindent{\it \textbf{Proof.}}
	\setlength{\parindent}{0pt}
	}
	\tcolorbox[blanker,breakable,left=5mm,parbox=false,
    before upper={\parindent15pt},
    after skip=10pt,
	borderline west={1mm}{0pt}{clmscol!50!\colormode}]
}{
    \textcolor{clmscol!50!\colormode}{\hbox{}\nobreak\hfill$\blacksquare$}
    \endtcolorbox
}

\NewDocumentCommand{\clmp}{m+m+m}{
    \begin{claim*}{#1}{}
        #2
    \end{claim*}

    \begin{clmpf}
        #3
    \end{clmpf}
}


% ============================
% Proof
% ============================

\newenvironment{pf}[1]{%
  \def\pfcolor{#1}% Store the color in a macro
  \noindent{\it \textbf{Proof.}}%
  \tcolorbox[blanker,breakable,left=5mm,parbox=false,%
    before upper={\parindent15pt},%
    after skip=10pt,%
    borderline west={1mm}{0pt}{\pfcolor},%
    coltext=\TextColor
  ]%
  \setlength{\parindent}{0pt}%
}{%
  \textcolor{\pfcolor}{\hbox{}\nobreak\hfill$\blacksquare$}%
  \endtcolorbox%
}

\newtcolorbox{solution}{
  blanker,
  breakable,
  left=5mm,
  parbox=false,%
  before upper={\parindent15pt},%
  after skip=10pt,%
  borderline west={1mm}{0pt}{\pbcolormain},%
  coltext=\TextColor
}


% ============================
% Remark
% ============================


\NewDocumentCommand{\rmk}{+m}{
    {\it \color{rmkscol!100!white}#1}
}

\newenvironment{remark}{
    \par
    \vspace{5pt}
    \begin{minipage}{\textwidth}
        {\par\noindent{\textbf{Remark.}}}
        \tcolorbox[blanker,breakable,left=5mm,
        before skip=10pt,after skip=10pt,
        borderline west={1mm}{0pt}{rmkscol!80!\colormode},
        coltext=\TextColor
        ]
}{
        \endtcolorbox
    \end{minipage}
    \vspace{5pt}
}

\NewDocumentCommand{\rmkb}{+m}{
    \begin{remark}
        #1
    \end{remark}
}


% Define a new command for problems
\renewcommand{\headrule}{\color{\TextColor}\hrulefill}
\newcommand{\C}{\mathbb{C}}
\newcommand{\R}{\mathbb{R}}
\newcommand{\Q}{\mathbb{Q}}
\newcommand{\Z}{\mathbb{Z}}
\newcommand{\N}{\mathbb{N}}
\renewcommand{\P}{\mathbb{P}}
\newcommand{\F}{\mathbb{F}}
\newcommand{\T}{\mathcal{T}}
\newcommand{\contr}{\Rightarrow \Leftarrow}
\newcommand{\s}{\(\,\)}
\renewcommand{\t}[1]{\text{#1}}
\newcommand{\ang}[1]{\left\langle #1 \right\rangle}
\newcommand{\coker}{\mathrm{coker}}

% Meta data
\setstretch{1.2}

\begin{document}
\color{\TextColor}
\pagecolor{\colormode}
\setlength{\parindent}{0pt}
\pagestyle{fancy}
\fancyhf{}
\fancyhead[LOE]{\color{\TextColor}\slshape{\textbf{Vincent Ferrara}}}
\fancyhead[ROE]{\color{\TextColor}\thepage}
\begin{problem}{}
    Let $X$ be a locally compact Hausdorff space. Prove that $X$ is
  second countable if and only if its one-point compactification is
  second countable.
\end{problem}
\begin{problemproof}
    $(\impliedby)$ Assume the one-point compactification of $X$ is second countable. Thus, it has a countable basis $\beta$.

    Since the topology of the one point compactification is made up of open sets of $X$ and sets of the form $\{X \cup \{\infty\}\} \setminus K$ s.t. $K$ is compact in $X$.

    Therefore, we know that for any $B \in \beta$ that $B \cap X$ is open. Since,

    if $B \subset X$ then it is already open in $X$.

    if not then $B=(X \cup \{\infty\}) \setminus K$. Thus, $B \cap X=X \setminus K$ which is open in $X$ since $K$ is compact thus closed. Thus, $X \setminus K$ is open.

    Thus, consider $\beta_X=\{B \cap X \mid B \in \beta\}$

    Let $x \in X$. Since $X$ is a subset of the one point compactification this means that $x \in B$ for some $B \in \beta$. Thus, $x \in B \cap X \in \beta_X$.

    If $x \in B_1 \cap B_2$ for some $B_1,B_2 \in \beta_X$. Then we know that $B_1=U_1 \cap X$ and $B_2=U_2 \cap X$ for $U_1,U_2 \in \beta$.

    Thus, $x \in U_1 \cap U_2$. Thus, there exists a $U_3 \in \beta$ s.t. $x \in U_3 \subset U_1 \cap U_2$.

    Thus, $x \in U_3 \cap X \subset B_1 \cap B_2$.

    Thus, $\beta_X$ is a countable basis for $X$.

    $(\implies)$ Assume $X$ is second countable.

    Let $\beta$ be a countable basis for $X$. 
    
    Define $\beta'=\{B \in \beta \mid \t{ there exists a compact } K_B \subset X \t{ with } B \subset K_B\}$.

    Let $x \in X$. Since $X$ is locally compact we know there is a neighborhood $U$ of $x$ s.t. $U \subset K$ for $K$ being compact.

    Since $\beta$ is a basis we know there is a $B$ s.t. $x \in B \subset U \subset K$. Thus, $\beta'$ covers $X$.

    Let $x \in B_1 \cap B_2$ for some $B_1,B_2 \in \beta'$. Thus, there is a $K_1,K_2$ s.t. these are compact and $B_1 \subset K_1$ and $B_2 \subset K_2$. Since $K_1,K_2$ are compact this implies they are closed. Thus, $K_1 \cap K_2$ is closed. Since $X$ is Hausdorff this means $K_1 \cap K_2$ is compact.

    Thus, $x \in B_1 \cap B_2 \subset K_1 \cap K_2$. Since $\beta$, we know there exists a $B_3 \in \beta$ s.t. $x \in B_3 \subset B_1 \cap B_2 \subset K_1 \cap K_2$ thus $B_3 \in \beta'$.

    Thus, $\beta'$ is a basis for $X$.

    Define $F=\{U \subset X \mid U=\bigcup_{i=1}^m B_i,\;\;B_i \in B',\;\;m \in \N\}$ which is countable since $B'$ is countable.

    Thus, for each set $U=B_1 \cup B_2 \dots \cup B_n \in F$ we know for each $B_i \subset K_i$ s.t. $K_i$ is compact. Thus, each $K_i$ is closed. Thus, the finite union of all the $K_i$'s is also closed, and since $X$ is Hausdorff this means the union is also compact. Call this union $K_U$.

    Let $\hat{X}$ be the one point compactification of $X$. 

    Define $V_U= \hat{X} \setminus K_U$

    Let $\hat{B}=\beta' \cup \{V_U \mid U \in F\}$

    Let $x \in \hat{X}$

    If $x \in X$ then $x \in B$ for some $B \in \beta'$

    If $x = \times$ then let $V$ be an open neighborhood of $\infty$. By definition of the topology on $\hat{X}$ we know that $V=\hat{X} \setminus K$ s.t. $K$ is compact.

    Since $\beta'$ is an open cover of $K$ we know there exists a finite subcover $B_1,\dots,B_n \in \beta$ s.t. $K \subset B_1 \cup \dots \cup B_n$.

    Thus, $K \subset K_1 \cup \dots \cup K_n=K_U$. Since each $B_i \subset K_i$ for some $K_i$ that is compact in $X$.

    Thus, $\infty \in V_U \subset V$.

    Let $U_1,U_2 \in \hat{B}$ s.t. $x \in U_1 \cap U_2$.

    If $x \in X$ then since open sets in $\hat{X}$ are either already open in $X$ or of the form $\hat{X} \setminus K$ where $K$ is compact. This means that the open sets intersected with $X$ are also open in $X$. We do not have to check for the first type but for the second type when intersecting with $X$ it becomes $X \setminus K$. Since $K$ is compact it is closed thus $X \setminus K$ is open.

    Therefore, any open set in $\hat{X}$ intersected with $X$ is open.

    Thus, $x \in U_1 \cap U_2 \cap X$. Since $\beta'$ is a basis this implies there exists a $B \in \beta'$ s.t. $x \in B \subset U_1 \cap U_2 \cap X$. Thus, $x \in B \subset U_1 \cap U_2$. Thus, done since $B \in \beta'$.

    If $x =\infty$ then $U_1 = \hat{X} \setminus K_{F_1}$ and $U_2 = \hat{X} \setminus K_{F_2}$. Thus, $K=K_{F_1} \cup K_{F_2}$ is also compact in $X$.

    Thus, by a similar proof above by showing $K$ is a subset of basis elements in $\beta'$ and taking a new set $K \subset K_U$. We know that $\infty \in V_U \subset \hat{X} \setminus K=U_1 \cap U_2$.

    Thus, $\hat{B}$ is a basis for $\hat{X}$ which is countable.
\end{problemproof}

\begin{problem}{}
    For each space below, determine which of the countability axioms
    (first countable, second countable, and separable) $X$ satisfies.
\end{problem}
\begin{subproblem}{}
    $X = \R^\omega$, equipped with the topology induced by the
    uniform metric.
\end{subproblem}
\begin{problemproof}
    Let $x \in \R^\omega$ with the uniform metric then consider $\{B(x,1/n) \mid n \in \N\}$.

    Let $U$ be an open neighborhood of $x$. Then, there exists an $\epsilon >0$ s.t. $B(x,\epsilon) \subset U$.

    Fix $N \in \N$ s.t. $1/N < \epsilon$. Thus, $B(x,1/N) \subset B(x,\epsilon) \subset U$.

    Thus, $\R^\omega$ is first countable.

    Assume $\R^\omega$ has a countable dense subset $U$.

    Consider $f: \{0,1\}^\omega \to U$ s.t. $f(x)=u$ s.t. $u \in B(x,1/2)$.

    Injective:

    Let $x,y \in \{0,1\}^\omega$ s.t. $f(x)=f(y)$. Thus, $x=y$ since $B(x,1/2)$ only contains one point from $\{0,1\}^\omega$.

    Thus, this implies $U$ is uncountable thus, $\R^\omega$ is not separable.

    Thus, since $\R^\omega$ is not separable proven in class this implies it is not second countable.
\end{problemproof}

\begin{subproblem}{}
    $X$ is the set $\R$ equipped with the topology generated by
    the basis of intervals of the form $[a, b)$, for $a, b \in \R$.
\end{subproblem}
\begin{problemproof}
    Let $x \in X$ then consider $\{[x,x+1/n) \mid n \in \N\}$.

    Let $U$ be an open neighborhood of $x$. Then, there exists an interval s.t. $x \in [a,b) \subset U$.

    Thus pick $N \in \N$ s.t. $a+1/N < b$. Thus, $x \in [a,a+1/N) \subset U$.

    Thus, $X$ is first countable.

    Consider $\Q$. $\Q$ is dense since any set $[a,b)$ with have a rational number in it. This is from $\Q$ being dense in $\R$ with standard topology. Thus, any open set intersects $\Q$ implying that $\overline{\Q}=\R$

    Thus, $X$ is separable.

    Assume for contradiction that $X$ is second countable. Thus, it has a second countable basis $\beta$.

    Define $f : X \to \beta$ s.t. $f(x)=B_x$ s.t. $x \in B_x \subset [x,x+1)$.

    Injective:

    Let $x,y \in X$ s.t. $f(x)=f(y)$. Thus, $B_x=B_y$. Thus, $B_y=B_x \subset [x,x+1)$ and $B_y=B_x \subset [y,y+1)$.

    Thus, $x=y$.

    This is a contradiction since $X$ is uncountable thus $X$ is not second countable.
\end{problemproof}

\begin{problem}{}
    Show that a compact metrizable space is always second
  countable. (This is exercise 4 in \S30 of Munkres; see that exercise
  if you need a hint.)
\end{problem}
\begin{problemproof}
    Let $X$ be compact and metrizable. Let $d$ be a metric on $X$.

    Consider $C=\{B(x,1/n) \mid x \in X\}$. This is an open cover of $X$. Since $X$ is compact we know there is a finite subcover $B_n$.

    Consider $\beta=\displaystyle\bigcup_{n \in \N}B_n$

    Let $x \in X$ since each $B_n$ is a finite subcover $x \in B \in B_n$.

    Let $x \in U_1 \cap U_2$ s.t. $U_1,U_2 \in \beta$. Thus, there is an $\epsilon > 0$ s.t. $x \in B(x,\epsilon) \subset U_1 \cap U_2$.

    Now consider $B_N$ s.t. $1/N < \epsilon/2$. Since $B_N$ is a finite subcover $x \in B$ for some $B=B(y,1/N) \in B_N$.

    Thus, for $z \in B$ we have $d(x,z) \leq d(x,y) + d(y,z) \leq 2/N < \epsilon$. Thus, $x \in B(y,1/N) \subset U_1 \cap U_2$.

    Thus, $\beta$ is a countable basis for $X$. Thus, $X$ is second countable.
\end{problemproof}

\begin{problem}{}
    Let $X_1, X_2$ be two spaces consisting of the same underlying set
  $X$ equipped with topologies $\T_1, \T_2$ respectively;
  suppose that $\T_1$ is finer than $\T_2$.
\end{problem}

\begin{subproblem}{}
    What (if anything) can you conclude about $X_1$ if $X_2$ is
assumed to be Hausdorff? What (if anything) can you conclude
about $X_2$ if $X_1$ is assumed to be Hausdorff?
\end{subproblem}
\begin{problemproof}
    Assume $X_2$ is Hausdorff. Let $x,y \in X_1$. Since $\T_2 \subset \T_1$. We know in $\T_2$ there exists disjoint nonempty open sets s.t. $x \in U$ and $y \in V$.

    Since $U,V \in \T_2 \subset \T_1$. Thus, $X_1$ is Hausdorff.

    The other way does not work. Consider $\T_1=\mathcal{P}(X)$, $\T_2=\{\emptyset,X\}$, and $X=\R$. $X_1$ is Hausdorff but $X_2$ is not.
\end{problemproof}

\begin{subproblem}{}
    Repeat part (a), but replace ``Hausdorff'' with ``regular''
and ``normal'' in turn.
\end{subproblem}
\begin{problemproof}
    Both ways do not work for regular. For one direction use the same counter example as above.

    For the other direction 
    
    Consider $\R=X_2$ with standard topology and $X_1=\R$ with a basis of $(a,b)$ and $(a,b) \setminus K$ s.t. $K=\{1/n \mid n \in \N\}$.

    Suppose there exists disjoint open sets containing $0,K$. Called $U,V$. Let $B$ be a basis element containing $0$ that is a subset of $U$. It must be a basis element of the form $(a,b) \setminus K$ since for any open interval containing zero intersects $K$.

    Fix $n \in \N$ large enough s.t. $1/n \in (a,b)$. Now pick a basis element about $1/n$ contained in $V$. It must be an interval $(c,d)$ since $1/n \in K$.

    This is a contradiction since fix $z \in \R$ s.t. $\max{c,1/(n+1)} < z < 1/n$. Thus, $z \in U \cap V$ which is a contradiction.

    Thus, $X_1$ is not regular but $X_2$ is.

    Both ways do not work for normal. For one direction use the same counter example as above.

    For the other direction

    We can also use the example of $\R$ with standard topology and $\R$ with the special $K$ topology. Since it is not regular it is also not normal, but $\R$ is normal.
\end{problemproof}

\begin{problem}{}
    Prove that if $M_1$ is an $n$-manifold and $M_2$ is an $m$-manifold,
  then the space $M_1 \times M_2$ is an $(n + m)$-manifold.
\end{problem}
\begin{problemproof}
    Let $U \times V$ be a basic open set in $M_1 \times M_2$. The reason we only have to look at basic open sets is because every set is a union of these so when applying a function $f(\bigcup B)=\bigcup f(B)$.

    Thus, $U$ is open in $M_1$ and $V$ is open in $M_2$.

    Thus, $U \cong S$ and $V \cong T$ s.t. $S$ is an open set in $\R^n$ and $T$ is an open set in $\R^m$.

    Let $f$ be a homeomorphism for $U$ to $S$ and let $g$ be the same for $V,T$.

    Consider, $h: U \times V \to S \times V$ s.t. $h(x,y)=(f(x),g(y))$. Thus, this is a homeomorphism since $f,g$ are homeomorphisms. Thus, $U \times V \cong S \times T \subset R^{n+m}$.

    Therefore, $M_1 \times M_2$ is a $n+m$ manifold.
\end{problemproof}

\begin{problem}{}
    Consider the cube $C = [-1,1] \times [-1,1] \times [-1,1]$. This
cube has six faces, given by the subsets
\begin{align*}
\{\pm 1\} \times [-1, 1] \times [-1, 1],\\
[-1, 1] \times \{\pm 1\} \times [-1, 1],\\
[-1, 1] \times [-1, 1] \times \{\pm 1\}.
\end{align*}
We can form the space $T_3$ by taking the quotient of $C$ by
``identifying opposite faces;'' precisely, by taking
$T_3 = C / \sim$, where $x \sim y$ if either $x = y$, or if any of
the following hold:
\begin{align*}
x = (\pm 1, u, v), \qquad y = (\mp 1, u, v)\\
x = (u, \pm 1, v), \qquad y = (u, \mp 1, v)\\
x = (u, v, \pm 1), \qquad y = (u, v, \mp 1)
\end{align*}
Prove that $T_3$ is homeomorphic to the product
$S^1 \times S^1 \times S^1$, and therefore a
manifold. \textit{Hint:} it suffices to find a continuous bijection
from $T_3$ to the product; why? You may also want to apply Problem 5
on Homework 4 (Theorem 22.2 in Munkres).
\end{problem}
\begin{problemproof}
    Let $p$ be the quotient map.

    Define $g : C \to S^1 \times S^1 \times S^1$ s.t. $g(x,y,z)=(e^{i\pi x},e^{i\pi x},e^{i\pi z})$.

    Since $e^{i \pi x}$ is continuous we know that $g$ is continuous.

    Let $x \neq y \in C$ s.t. $x \sim y$.

    Thus, let $x=(\pm 1,u,v)$ and $y=(\mp 1, u ,v)$. However, since $e^{-i\pi}=e^{i\pi}$. This, implies that $g(x)=g(y)$. The same argument works for the other types of equivalences.

    Thus, $g$ is constant on the equivalence classes. Therefore, by Theroem 22.2 we get an induced map $f : T_3 \to S^1 \times S^1 \times S^1$ s.t. $f \circ p = g$.

    Also, since $g$ is continuous we know that $f$ is continuous.

    Injective:

    Let $a,b \in T_3$ s.t. $f(a)=f(b)$.

    Since $q$ is surjective this implies there is $x,y \in C$ s.t. $q(x)=a$ and $q(y)=b$.

    Thus, $f(q(x))=f(q(b))$. Therefore, $g(x)=g(y)$.

    Thus, $(e^{i\pi x_1},e^{i\pi x_2},e^{i\pi x_3})=(e^{i\pi y_1},e^{i\pi y_2},e^{i\pi y_3})$

    Thus, $e^{i \pi (x_j - y_j)}=1$

    This, implies that $\pi(x_j - y_j)=2 \pi k$ for some $k \in \Z$.

    Dividing by $\pi$ gets:

    $x_j-y_j = 2k$

    Since $x_j,y_j \in [-1,1]$. This implies that $k$ is 0, 1, or $-1$.

    Therefore this implies that either:
    \begin{align*}
        x_j-y_j=0 \implies x_j=y_j\\
        x_j-y_j=2 \implies x_j=1, \; y_j=-1\\
        x_j-y_j=-2 \implies x_j=-1, \; y_j=1
    \end{align*}

    Thus, this implies that $x \sim y$.

    Therefore, $a=b$. Thus, $f$ is injective.

    Surjective:

    Let $x \in S^1 \times S^1 \times S^1$. Therefore, since $e^{i\pi t} : [-1,1] \to S^1$ is surjective there is $x_1,x_2,x_3$ s.t. $x=(e^{i\pi x_1},e^{i \pi x_2}, e^{i \pi x_3})$.

    Thus, consider $f(q(x_1,x_2,x_3))=g(x_1,x_2,x_3)=x$.

    Therefore, $f$ is surjecitve.

    First, since $C$ is closed and bounded we know that it is compact. Thus, $q(C)=T_3$ is compact. Also, since $S_1 \times S_1 \times S_1$ is a subspace of a Hausdorff space it is Hausdorff. Thus, since $f$ is a continuous bijection form a compact space to a Hausdorff space it is a homeomorphism. Therefore, $T_3 \cong S^1 \times S^1 \times S^1$.
\end{problemproof}
\end{document}